%!TEX root = main.tex

\section{Introduction}
\label{sec:introduction}
\IEEEPARstart{R}{obust}
% Robust 
estimation is a crucial tool for robotics and computer vision, 
%and is 
being
concerned with the estimation of unknown quantities 
(\eg the state of a robot, or of variables describing the external world) 
from noisy and potentially corrupted measurements. Corrupted measurements (\ie \emph{outliers}) can be caused by sensor malfunction, but are  more commonly associated with % caused by
incorrect data association \mbox{and model misspecification~\cite{Cadena16tro-SLAMsurvey,Yang19rss-teaser}.} 

%%%%%%%%%%%%%%%%%%%%%%%%%%%%%%%%%%%%%%%%%%%%%%%%%%%%%%%%%%%%%%%%%%%%%%%%%%%%%%%%%%%%%%%%%%%%%%%%%%%%%%%%%%%%
%% ESTIMATION WITHOUT OUTLIERS
In the outlier-free case, common estimation problems 
%in 
% Many estimation problems in 
% robotics and vision 
are formulated as a least squares optimization: % in the form:
%terms of the following general form of optimization:
% seek to solve 
% the following geneal form of optimization:
\bea \label{eq:leastSquares}
\min_{\vxx \in \calX} \;\sumAllPointsi r^2(\vy_i,\vxx),
\eea 
where $\vxx$ is the variable we want to estimate (\eg the pose of an unknown object); 
$\calX$ is the domain of $\vxx$ (\eg the set of 3D poses);
$\vy_i$ ($i=1,\ldots,N$) are given measurements (\eg pixel observations of points belonging to the object); 
  and the function $r(\vy_i,\vxx)$ is the \emph{residual} error for the $i$-th measurement,
quantifying the mismatch between the expected measurement at an estimate $\vxx$ and the actual measurement~$\vy_i$.
% The mathematical expression of $r(\vy_i,\vxx)$ varies from problem to problem, 
% where $\calK\subseteq\Real{n}$ is the (potentially non-convex,~\eg~$\SOthree$) feasible set, $r(\calD_i,\vxx) := r_i \in \Real{}$ is the $i$-th \emph{residual} computed by applying the residual function $r(\cdot,\cdot)$ on the $i$-th measurement data $\calD_i$ and the decision variable $\vxx$. When there are zero noise among all the measurements $\calD_i$'s, then $r_i=0$ for all $i=1,\dots,N$.
In typical robotics and computer vision applications, 
\revise{~the least squares optimization}~\eqref{eq:leastSquares} is difficult to solve globally, due to the nonlinearity of the 
residual errors and the nonconvexity of the domain~$\calX$. Despite these challenges, 
the research community has developed closed-form solutions and\revise{~globally optimal} solvers for many such problems.
% of interest.
Specifically, 
while closed-form solutions are rare~\cite{Horn87josa}, 
  %in recent years researchers have %developed global solvers shown that 
  % researchers 
\emph{Semidefinite Programming} (SDP) and \emph{Sums-of-Squares} (SOS)~\cite{Blekherman12Book-sdpandConvexAlgebraicGeometry}
%~\cite{Lasserre07siam-sos,Parrilo00thesis} 
relaxations have been recently shown to 
  be a powerful tool to obtain \emph{certifiably optimal} solutions to relevant instances of problem~\eqref{eq:leastSquares}, 
  ranging from pose graph optimization~\cite{Carlone16tro-duality2D,Rosen18ijrr-sesync},
   rotation averaging~\cite{Eriksson18cvpr-strongDuality},
   anisotropic registration~\cite{Briales17cvpr-registration}, 
   two-view geometry~\cite{Briales18cvpr-global2view}, 
   and PnP~\cite{Agostinho2019arXiv-cvxpnpl}.
   The resulting techniques are commonly referred to as \emph{non-minimal solvers}, to contrast them against 
   methods that solve problem~\eqref{eq:leastSquares} using only a small (minimal) subset of measurements $\vy_i$ 
   (see related work in Section~\ref{sec:relatedWork}).


%%%%%%%%%%%%%%%%%%%%%%%%%%%%%%%%%%%%%%%%%%%%%%%%%%%%%%%%%%%%%%%%%%%%%%%%%%%%%%%%%%%%%%%%%%%%%%%%%%%%%%%%%%%%
%% ESTIMATION WITH OUTLIERS
%   $\rho(\cdot)$ is a given cost function.
% In eq.~\eqref{eq:generalEstimation}, 
Unfortunately, in the presence of outliers, problem~\eqref{eq:leastSquares}'s solution provides a 
poor estimate for 
%the variable 
$\vxx$. This limits the applicability of existing non-minimal solvers, 
% Indeed, the quadratic cost typically follows from the assumption that
% the measurement noise is Gaussian~\cite{Cadena16tro-SLAMsurvey} and is less suitable to deal with heavy-tailed noise. 
 since they can be applied only after the outliers have been removed. 
 Yet,
 the theory of robust estimation  
 suggests regaining robustness by substituting the quadratic cost in\revise{~the least squares problem}~\eqref{eq:leastSquares} 
 with a \emph{robust cost} $\rho(\cdot)$:
%
\bea \label{eq:generalEstimation}
\min_{\vxx \in \calX} \; \sumAllPointsi \rho( r(\vy_i,\vxx) ).
\eea 
%
% where $\vxx$ is the variable we would like to estimate (\eg the pose of an unknown object), 
% $\calX$ is the domain of $\vxx$ (\eg the set of 3D poses),
% $\vy_i$ ($i=1,\ldots,m$) are given measurements (\eg pixel observations of points belonging to the object), 
%   and $\rho(\cdot)$ is a given cost function.
% In eq.~\eqref{eq:generalEstimation}, the function $r(\vy_i,\vxx)$ is the \emph{residual} error for the $i$-th measurement,
%  that quantifies the mismatch between the expected measurement given an estimate $\vxx$ and the actual measurement $\vy_i$.
For instance, $\rho(\cdot)$ can be a Huber loss, a truncated least squares cost, or a Geman-McClure cost~\cite{Black96ijcv-unification}.
% For instance, $\rho(\cdot)$ can be a Huber loss, a truncated least squares cost, or a Geman-McClure cost~\cite{Black96ijcv-unification}.  
To date, the application of non-minimal solvers  to eq.~\eqref{eq:generalEstimation} has been limited. 
Some of the non-minimal solvers designed for~\eqref{eq:leastSquares} can be extended to include a \emph{convex} robust cost function, such as the Huber loss~\cite{Carlone18ral-robustPGO2D}; however, it is known that convex losses have a low breakdown point and are still sensitive to gross outliers~\cite{Lajoie19ral-DCGM}. 
In rare cases, the literature provides \emph{robust non-minimal solvers} that achieve 
global solutions to specific instances of problem~\eqref{eq:generalEstimation}, 
such as robust registration~\cite{Yang19rss-teaser,Yang19iccv-QUASAR}. 
%, and pose graph optimization~\cite{Lajoie19ral-DCGM}.  
However, 
these solvers cannot be easily extended to other estimation problems, and they rely on solving large SDPs, which is currently
 % might be
  impractical for large-scale problems.
% Currently, however, the application of eq.~\eqref{eq:generalEstimation}'s approach to non-minimal solvers in robotics and computer vision has been limited. 
% Particularly, 
% although
% some of the non-minimal solvers above can be indeed extended to include a \emph{convex} robust cost (such as the Huber loss~\cite{Carlone18ral-robustPGO2D}),
% %. But 
% such costs have low breakdown points, and are still sensitive to gross outliers~\cite{Lajoie19ral-DCGM}. 
% And although for special cases the literature provides \emph{robust non-minimal solvers} achieving
% global solutions to specific instances of problem~\eqref{eq:generalEstimation}
% (such as robust registration~\cite{Yang19rss-teaser,Yang2019ICCV-QUASARforWahba}, and pose graph optimization~\cite{Lajoie19ral-DCGM}),
% %.  However, 
% these solvers cannot be easily extended to other estimation problems, and rely on solving large SDPs, which 
% \mbox{might be impractical for very large-scale problems.}
 %real-time
% is impractical 
% require specialized solvers~\cite{Rosen18ijrr-sesync}.
%which is practical only in relatively small instances. %\red{HY: should mention BnB here as well?} 
% Alternative global approaches to outlier rejection are based on branch-and-bound (\bnb)~\cite{Bustos18pami-GORE,chin2017slcv-maximumconsensusadvances}, which unfortunately
%  has exponential runtime in the worst case.
% we provide a specialized global solver : Teaser/Quasar- but there are few of them and there are no scalable implementations
% others are exponential time.


{\bf Contributions.} In this paper, we aim to reconcile non-minimal solvers and robust estimation, by providing a 
\textit{general-purpose}\footnote{For a given problem~\eqref{eq:generalEstimation}, we only assume the existence of a non-minimal solver for the 
corresponding outlier-free problem~\eqref{eq:leastSquares}.} algorithm to solve problem~\eqref{eq:generalEstimation} without requiring an initial guess. 
 % Our only assumption is that a non-minimal solver exists 
% , for any instance 
% without requiring an initial guess; 
%  our only restriction is to problems 
% our algorithm only requires the existence of a non-minimal solver for the 
%  We only require 
% for any instance where a non-minimal solver is available, and without an initial guess. 
We achieve this by combining  non-minimal solvers with an approach known as \emph{graduated non-convexity} (\GNC). 
% This paper aims to develop general-purpose algorithms for a broad class of estimation problems 
In contrast, standard algorithms for problem~\eqref{eq:generalEstimation} rely on iterative optimization to refine a given initial guess, which causes the result to be brittle when the quality of the  guess is poor. 
Particularly, we propose three contributions.

First, we revisit the Black-Rangarajan duality~\cite{Black96ijcv-unification} between robust estimation and outlier processes.
%, and apply it to spatial perception problems in robotics.
 % In particular, we focus on 
 % The third contribution is to demonstrate how two examples of the \GNC framework, the \GNC-\emph{Geman McClure} (\GNCGM) and the \GNC-\emph{Truncated Least Squares} (\GNCTLS), can be applied on four robotics and computer vision applications: point cloud registration, generalized 3D registration, pose graph optimization, and shape alignment from a single image. 
  We also revisit the use of \emph{graduated non-convexity} (\GNC) as a general tool to solve a non-convex 
  optimization problem without an initial guess.
  While Black-Rangarajan duality and \GNC have been used in early vision problems, such as stereo reconstruction, 
  image restoration and segmentation, and optical flow, 
%    beautiful theory of Black-Rangarajan duality has been confined to simple estimation problems
% mainly for early vision problems, including stereo reconstruction, image restoration and segmentation, and optical flow
% -- mainly focuses on surface fitting
 we show that combining them with non-minimal solvers allows solving \emph{spatial perception} problems, ranging from 
 mesh registration, pose graph optimization, and image-based object pose estimation.
 %it provides a general tool for robust estimation in many spatial 
 % The crux is that while the original theory, 

Our second contribution is to
tailor Black-Rangarajan duality and \GNC to the 
 \textit{Geman-McClure} and \textit{truncated least squares} costs.
We show how to optimize these functions  by alternating two steps: 
a \emph{variable update}, which solves a \emph{weighted} least squares problem using non-minimal solvers; 
 and a \emph{weight update}, \mbox{which updates the outlier process in closed~form.} 
 % In addition, we show that the use of \emph{graduated non-convexity} (\GNC) 
 %framework, where a family of optimizations with increasing non-convexity are solved to gradually reject outliers. 
 %  In addition, lack of efficient \emph{minimal} solvers prevents us from applying \ransac to gain robustness. In this paper, our first contribution is to propose the \emph{Graduated Non-Convexity} (\GNC) framework, where a family of optimizations with increasing non-convexity are solved to gradually reject outliers. 
 % Our second contribution is to show that solving each optimization in the \GNC framework can be done in two steps: 
% Our second contribution 

Our approach requires 
%the availability of 
a non-minimal solver, but currently there is no such solver for 
image-based object pose estimation (also known as \emph{shape alignment}).
%Therefore, 
Our third contribution is to present a novel non-minimal solver for shape alignment.
While related techniques propose approximate relaxations~\cite{Zhou17pami-shapeEstimationConvex}, we 
provide a 
\emph{certifiably optimal} solution 
%can be obtained by applying 
using SOS relaxation. 

We demonstrate our \emph{robust non-minimal solvers} on point cloud registration (\PCR), mesh registration (as known as \emph{generalized} registration, \GReg), pose graph optimization (\PGO), and shape alignment (\shape).
%point cloud and mesh registration, pose graph optimization, and image-based object pose estimation.
 % Replacing the least squares cost with robust cost functions yields non-convex optimizations whose performance are highly sensitive to initial guesses. 
Our solvers  are robust to {70-80\%} outliers, 
%across applications, 
outperform \ransac, 
are more accurate than specialized \emph{local} solvers, and 
faster than \emph{global} solvers.
% The proposed solvers are fast and general-purpose, are robust to at least \red{70-80\%} of outliers, are superior and more broadly applicable than existing general-purpose methods (\eg{ }\ransac), \red{and perform {comparably to specialized methods.}}
%, while being 
%fast and  general-purpose. 
% The paper also extends the literature on non-minimal solvers by proposing an SOS solver for image-based object pose estimation.
% issue with BR: nonconvex domain.. even with continuation, problems is non-convex
% while these techniques are not necessarily as robust as specialized methods, they carry the promise of fast 
% execution and practical robustness
% faster than specialized methods
%  beautiful theory of Black-Rangarajan duality has been confined to simple estimation problems
% mainly for early vision problems, including stereo reconstruction, image restoration and segmentation, and optical flow
% % -- mainly focuses on surface fitting
% -- later for:" A graduated assignment algorithm for graph matching"
% -- markov random fields
% The main contribution of this paper is to bridge the gap between non-minimal solvers and robust estimation, 
%  and provide a general-purpose tool for robust global estimation, which can be applied to any problem where a non-minimal solver is available.
% minimal solvers
% 1) the chapter is about GNC but NOT robust estimation (only about early vision), but has alternation
% 2) Black-Rangarajan duality is about line processes, mentions continuation = graduated nonconvexity and alternation
% 3) FRG has both line processes, GNC, and alternation but focuses on a specific application
 % Towards this goal, we leverage the Black-Rangarajan duality between robust estimation and outlier processes 
 % (which has been traditionally applied to early vision problems) and show that the 
 % joint use of non-minimal solvers and \emph{graduated non-convexity} (\GNC) allows computing robust solutions without requiring 
 % an initial guess. % for the variables. 
 % 

